\section{Limitations}

\begin{itemize}
  \item \textbf{Local students see frames, not video.} The student receives up to 4 frames per second at 512 pixels
  (about 70 tokens a frame), without audio. Code2Video's own protocol (Gemini on the mp4) is implemented and exported
  for grading by hand, but not run.
  \item \textbf{Few students.} The main student is \texttt{gemma4:26b}; the prior was also measured with
  \texttt{qwen2.5vl:7b}. On NOVEL-12 we also asked Claude Haiku 5.5 and Sonnet 5.5, and they agree with
  \texttt{gemma4:26b} on the three findings we tested: the videos teach, the student learns the video's own topic,
  and the Claude fixes give no gain beyond noise. But Claude also wrote the questions and made the layout fixes, and it returns
  no log-probabilities, so its answers are scored right or wrong. Open students of other families are untested.
  \item \textbf{Invented concepts are small.} NOVEL-12's concepts are definable on one card; they test whether a video
  transmits a definition and its consequences, not deep understanding of a large theory. TEST was written by one model
  (Claude) and not reviewed by teachers.
  \item \textbf{The evidence judge is not calibrated.} Its decisions have not been compared with human annotations
  (e.g. Cohen's $\kappa$~\cite{cohen1960}); the forced-choice versions do not use it.
  \item \textbf{Twelve topics per set.} Means over 12 topics move a lot with one topic. We do not pool questions
  across topics, because the topic, not the question, is what receives a repair; a larger set would make the means
  steadier.
  \item \textbf{The local coder.} After one retry, two of twelve NOVEL-12 videos were still incomplete and only
  \AdmClean{} of \AdmTotal{} were clean. Stronger code models are future work; the Claude-fixed versions show the
  room.
  \item \textbf{The repair study is small.} All six stages are scored, on 10 concepts with one student. From F1 on,
  the gate kept the learner revision on at most half of them. A gain beyond noise across concepts was therefore out of
  reach, and small benefits of the repair remain possible. The two measures also disagree: kept revisions raise the
  gain on 3 frames (M9) but not on the dense video (M7).
  \item \textbf{The model-free checks are narrow.} The timing check covers two kinds of fault; the attachment check
  skips labels anchored to a text glyph.
  % Hand-typed from bench/layout_study.log, layout_study_rerun.log, layout_study_attach.log (sum of 'conflicts'), and
  % from bench/results/layout-study/rows.json of 22:33 (fixed tracker, commit 937336d), summed over the 10 admitted concepts.
  \item \textbf{The conflict counts were measured with a tracker bug.} The tracker sometimes gave a deleted object's
  code line to a new one, so the same code gave 29, 29 and 31 lasting conflicts on three runs. With the bug fixed,
  the counts repeat exactly. Our videos then have 27 lasting conflicts, as do our videos with timing fixes, against 4
  for the Claude-fixed versions. The tables keep the earlier counts. MVQS, the attachment counts and every
  conclusion are unchanged.
  \item \textbf{Counterfactual relabelling was not applied to MMMC.} It would rewrite TEST per video and break the
  paired design; NOVEL-12 controls the prior instead.
\end{itemize}

\section{Future work}

\begin{itemize}
  \item \textbf{A gate that sees what a viewer sees.} The repair gate checks only the layout score, and it let through
  revisions that added conflicts and broken attachments. A gate that also counts these faults is built (crowding,
  conflicts, attachments) but not yet evaluated.
  \item \textbf{A larger study.} With 10 concepts of 8 questions, only effects of about 24 points or more are
  detectable. More concepts, and a gate that keeps more revisions, would let smaller effects of layout and repair show.
  \item \textbf{Stronger coders and a fixed layout.} A study with Claude models as planner and coder, and with
  Code2Video's anchor-grid layout, was running at submission. It asks whether better code, rather than repair,
  is what makes videos teach more.
  \item \textbf{Students that watch video.} Code2Video's protocol (a model reading the mp4 with audio) is implemented
  but not run; open students of other families are untested.
  \item \textbf{A layout score that measures reading.} MVQS measures box separation. A score built on attachments, on
  what is on screen when, and on crowded text would track what a viewer can read.
  \item \textbf{Cheaper scoring.} The dense video is about two thirds of the student's work. Scoring on 3 frames per
  section only (M9) is implemented as an option, for studies that do not need M7.
\end{itemize}

\section{Conclusion}

TeachQuiz, as published and as we first ran it, could not have detected teaching: a grading bug dropped answers,
the student's refusals set $S_1$, and the questions were answered from prior knowledge. Forced choice, a measured
prior, normalisation, dense silent frames, a motion control, continuous scoring and option rotation each address a
named failure, and the ladder shows which of them a teaching metric actually needs. A benchmark whose prior is
controlled by construction is the precondition for asking whether a repair teaches. The layout score has the same
weakness from the other side: it measures box separation, not whether a viewer can read what belongs to what. On
NOVEL-12, the repair loop gives no gain beyond noise at any stage. None of the added features makes the student's
answers beat a repair of the same size made without them.
